%%% ---------------------------------------------------------------------------
%%% main.tex —— 主控文件
%%%
%%% 这个文件只负责「装配」，不放正文、不放格式。
%%%   要改样式  → acadarticle.cls
%%%   要改元数据 → setup.tex
%%%   要加宏包/自定义命令 → preamble.tex
%%%   要写正文  → sections/*.tex
%%%
%%% 编译：make        （或 latexmk main.tex）
%%% ---------------------------------------------------------------------------
\documentclass[
  lang      = zh,        % zh | en
  fontset   = fandol,    % fandol（零配置） | noto（更好看，需 noto-fonts-cjk） | none
  mathstyle = iso,       % iso（遵循 GB/T 3101） | tex
  bibstyle  = gb7714,    % gb7714 | gb7714ay | ieee | numeric | authoryear | none
  % anonymous,           % 盲审：隐去作者/单位/基金/致谢
  % linenumbers,         % 送审稿加行号
  % print,               % 纸质版：超链接转黑色
  % twocolumn,           % 双栏（原样传给 article）
]{acadarticle}

%%% ---------------------------------------------------------------------------
%%% setup.tex —— 论文元数据
%%%
%%% 只填内容，不写格式。所有排版规则在 acadarticle.cls 里。
%%% 对应 GB/T 7713.2—2022 前置部分的著录项。
%%% ---------------------------------------------------------------------------

%%% ---- 中文题名与作者 ----
%%% 题名一般不超过 20 字，不使用非公知的缩略语、字符、代号
\title{基于多尺度特征融合的遥感图像目标检测方法}

%%% 多作者用 \quad 分隔；上标数字对应下面的单位编号
\author{张三\textsuperscript{1}\quad 李四\textsuperscript{1,2}\quad 王五\textsuperscript{2}}

%%% 单位要写全称，到二级单位，附城市与邮编（GB/T 7713.2）
\affiliation{%
  1.~中国科学技术大学 信息科学技术学院，合肥 230027；\\
  2.~中国科学技术大学 大数据学院，合肥 230026%
}

%%% ---- 关键词：3~8 个，分号分隔，与英文一一对应 ----
\keywords{遥感图像；目标检测；多尺度特征；特征融合；卷积神经网络}

%%% ---- 中图分类号与文献标志码 ----
%%% 中图分类号查《中国图书馆分类法》；文献标志码 A=理论与应用研究论文
\clcnumber{TP391.4}
\doccode{A}

%%% ---- 英文题名、作者、单位（对应中文，供英文摘要块使用）----
\entitle{Multi-scale Feature Fusion for Object Detection in Remote Sensing Imagery}
\enauthor{ZHANG San\textsuperscript{1}, \quad LI Si\textsuperscript{1,2}, \quad WANG Wu\textsuperscript{2}}
\enaffiliation{%
  1.~School of Information Science and Technology, University of Science and
  Technology of China, Hefei 230027, China;\\
  2.~School of Data Science, University of Science and Technology of China,
  Hefei 230026, China%
}
\enkeywords{remote sensing imagery; object detection; multi-scale feature;
  feature fusion; convolutional neural network}

%%% ---- 基金项目与通信作者（自动排为首页无标记脚注；anonymous 模式下隐去）----
\funding{国家自然科学基金资助项目（62376001）；国家重点研发计划（2023YFB0000000）}
\corresponding{李四，教授，E-mail: lisi@ustc.edu.cn}

%%% ---- 页眉的短标题 ----
\runningtitle{基于多尺度特征融合的遥感图像目标检测方法}

%%% ---- 日期：留空则不显示 ----
\date{}

%%% ---- 参考文献数据库 ----
%%% refs.bib 是指向 ../common/refs.bib 的符号链接，与 beamer 模板共用同一份
%%% 文献库，改一处两处都生效。
\addbibresource{refs.bib}
       % 题名、作者、单位、关键词、基金……
%%% ---------------------------------------------------------------------------
%%% preamble.tex —— 你自己的宏包与自定义命令
%%%
%%% 类文件里已经载入了：
%%%   ctex / fontspec / geometry / unicode-math / acadmath / amsmath / mathtools
%%%   graphicx / booktabs / array / multirow / makecell / longtable / threeparttable
%%%   caption / subcaption / float / algorithm / algpseudocode / amsthm
%%%   biblatex / fancyhdr / xcolor / hyperref / bookmark / cleveref
%%%   enumitem / url / microtype
%%% 不要重复载入这些，尤其不要重新载入 hyperref。
%%%
%%% 共享的数学记号（\vect \mat \norm \argmin \dd \subtext …）见 ../common/acadmath.sty，
%%% 那个文件同时被 beamer 模板使用，改它可以让论文和报告的符号保持一致。
%%% ---------------------------------------------------------------------------

%%% ---- 按需追加的宏包 ----
% \usepackage{tikz}
% \usepackage{pgfplots}\pgfplotsset{compat=1.18}
% \usepackage{listings}          % 代码
% \usepackage{physics}           % 注意：与 unicode-math 有已知冲突，慎用

%%% ---- 本文专用的记号 ----
% 例：把常用的模型名、数据集名定成命令，避免全文写法不一致
\newcommand{\ourmethod}{MSFF-Net}
\newcommand{\dataset}{DOTA-v2.0}

% 例：本文专用的数学量
\newcommand{\featmap}{\mat{F}}          % 特征图
\newcommand{\loss}{\mathcal{L}}         % 损失函数
\newcommand{\params}{\vect{\theta}}     % 参数向量

%%% ---- 占位图 ----
%%% 只是为了让模板在没有任何图片文件时也能编译。
%%% 真正插图时删掉它，改用：\includegraphics[width=0.85\linewidth]{architecture}
%%% 用法：\placeholder{<宽>}{<高>}{<说明文字>}
\newcommand{\placeholder}[3]{%
  \begingroup
    \fboxsep=0pt \fboxrule=0.4pt
    \fbox{\parbox[c][#2][c]{\dimexpr#1-2\fboxrule\relax}{\centering\footnotesize\ttfamily\color{gray}#3}}%
  \endgroup
}

%%% ---- 若要覆盖类文件里的样式，写在这里（会后于类文件生效）----
% 例：改正文字号
% \AtBeginDocument{\zihao{-4}}
%
% 例：图题改成左对齐
% \captionsetup[figure]{justification=raggedright, singlelinecheck=false}
    % 你自己的宏包与命令

\begin{document}

%%% ---- 前置部分（GB/T 7713.2 规定顺序：题名→作者→单位→摘要→关键词）----
\maketitle
% !TeX root = ../main.tex
%%% ---------------------------------------------------------------------------
%%% 摘要与关键词
%%%
%%% GB/T 7713.2 对摘要的要求（写之前先看一遍）：
%%%   · 报道性摘要，能脱离全文独立存在
%%%   · 不出现「本文」「作者」等自指表述
%%%   · 不引用文献编号、不出现图表号
%%%   · 不使用非公知的缩略语
%%%   · 中文摘要一般 300 字以内，英文摘要与中文对应
%%%   结构建议：目的 → 方法 → 结果（带具体数据）→ 结论
%%%
%%% 关键词写在 setup.tex 里，这里的环境会自动排出来。
%%% ---------------------------------------------------------------------------

\begin{abstract}
针对遥感图像中目标尺度跨度大、小目标特征在深层网络中易丢失的问题，提出一种
多尺度特征融合的目标检测方法。该方法在特征金字塔的基础上引入跨层注意力
门控，按空间位置自适应地加权来自不同下采样倍率的特征，并以可变形卷积
对齐相邻层级的感受野偏移。在 \dataset{} 数据集上的实验表明，所提方法的
平均精度均值达到 \qty{78.4}{\percent}，较基线方法提升 \qty{3.7}{\percent}，
其中小目标类别的平均精度提升 \qty{6.2}{\percent}，推理速度保持在
\qty{31}{\fps}。结果说明跨层注意力门控能够在不显著增加计算量的
前提下缓解小目标特征衰减，为高分辨率遥感图像的实时检测提供了可行方案。
\end{abstract}

%%% 英文摘要块（lang=en 时删掉这一段即可）
\begin{enabstract}
To address the large scale variation of objects and the loss of small-object
features in deep networks for remote sensing imagery, a multi-scale feature
fusion detection method is proposed. Built upon a feature pyramid, the method
introduces a cross-level attention gate that adaptively weights features from
different downsampling rates at each spatial location, and employs deformable
convolution to align the receptive-field offsets between adjacent levels.
Experiments on \dataset{} show that the proposed method attains a mean average
precision of \qty{78.4}{\percent}, an improvement of \qty{3.7}{\percent} over
the baseline, with a \qty{6.2}{\percent} gain on small-object categories while
maintaining an inference speed of \qty{31}{\fps}. The results
indicate that the cross-level attention gate mitigates small-object feature
attenuation without a significant increase in computation, offering a practical
route to real-time detection in high-resolution remote sensing imagery.
\end{enabstract}


%%% ---- 主体部分 ----
% !TeX root = ../main.tex
%%% ---------------------------------------------------------------------------
%%% 引言
%%%
%%% 结构建议（学术写作通行的「漏斗式」）：
%%%   1. 研究领域的重要性与应用背景
%%%   2. 已有工作及其成果
%%%   3. 已有工作的不足 —— 这是全文的立论点，要具体，不要泛泛说「效果不佳」
%%%   4. 本文的思路与主要贡献（分条列出）
%%%   5. 全文组织
%%%
%%% 引用一律用 \cite{...}；不要手写 [1]。
%%% ---------------------------------------------------------------------------
\section{引言}
\label{sec:intro}

高分辨率遥感图像的目标检测是对地观测、灾害评估与城市管理的关键技术
环节\cite{cheng2016survey}。与自然图像相比，遥感图像的目标具有尺度跨度大、
方向任意、背景杂乱三个显著特点，其中尺度问题尤为突出：同一景中大型建筑
可占据数千像素，而车辆、船舶等目标往往只有十余像素。

围绕这一问题，已有工作主要沿两条路线展开。一是构建多尺度特征表示，
以特征金字塔网络\cite{lin2017fpn}为代表，通过自顶向下的横向连接把深层
语义信息回传到浅层；二是改进检测头与标签分配策略\cite{zhang2020atss}，
使不同尺度的目标获得更合理的正样本。国内学者在此基础上针对遥感场景
提出了旋转框表示与角度回归的改进\cite{wang2021remote}。

然而，现有方法在处理小目标时仍存在两点不足：

\begin{enumerate}
  \item 特征金字塔的横向连接采用固定权重相加，未考虑不同空间位置对各
        层级特征的需求差异——背景区域与小目标区域被同等对待；
  \item 相邻层级之间存在因下采样引入的感受野偏移，直接逐元素相加会
        造成特征错位，这一误差在小目标上被放大。
\end{enumerate}

针对上述问题，本文提出\ourmethod{}。主要贡献如下：

\begin{enumerate}
  \item 提出跨层注意力门控，按空间位置自适应地加权各层级特征，使小目标
        区域更多地保留浅层高分辨率信息；
  \item 引入可变形卷积对齐相邻层级的感受野偏移，缓解逐元素融合造成的
        特征错位；
  \item 在 \dataset{} 上验证了方法的有效性，并通过消融实验分离了两个
        模块各自的贡献。
\end{enumerate}

全文组织如下：\cref{sec:related} 综述相关工作；\cref{sec:method} 给出
\ourmethod{} 的结构与理论分析；\cref{sec:exp} 报告实验结果；
\cref{sec:conclusion} 总结全文。

% !TeX root = ../main.tex
%%% ---------------------------------------------------------------------------
%%% 相关工作
%%%
%%% 综述不是文献罗列。每一段应当围绕**一条技术脉络**，说清楚
%%% 「大家怎么做的 → 共同的局限 → 与本文的关系」，最后落回本文的切入点。
%%% ---------------------------------------------------------------------------
\section{相关工作}
\label{sec:related}

\subsection{多尺度特征表示}
\label{sec:related:multiscale}

图像金字塔是最早的多尺度方案，代价是重复的前向计算。特征金字塔
网络\cite{lin2017fpn}把多尺度构造转移到特征层面，以近乎不变的计算量获得
了多层级表示，此后成为检测器的标准组件。后续工作从连接拓扑
入手做了大量改进，但横向连接的融合权重在空间上仍是均匀的。

\subsection{注意力机制在检测中的应用}
\label{sec:related:attention}

通道注意力与空间注意力已被广泛用于强化判别性特征。在检测任务中，
注意力多被放在骨干网络内部，用于抑制背景响应；将其用于\emph{层级选择}
——即决定某个空间位置应当更多依赖哪一层级的特征——的研究相对较少，
这正是本文的切入点。

\subsection{遥感图像目标检测}
\label{sec:related:rsdet}

遥感场景的特殊性催生了旋转框表示、角度周期性损失以及大幅面图像的
切片推理策略\cite{wang2021remote}。这些工作与本文正交：本文改进的是
特征融合方式，可以直接嵌入到上述检测框架中。

% !TeX root = ../main.tex
%%% ---------------------------------------------------------------------------
%%% 方法
%%%
%%% 这一节演示了公式、定理、算法的规范写法，改写时请保留这些约定：
%%%   · 变量斜体、单位正体、函数名正体（GB/T 3101，见 ../../STANDARDS.md §1.3）
%%%   · 微分写 \dd x，自然常数写 \e，虚数单位写 \ii
%%%   · 说明性下标写 \subtext{max}，不要写 _{max}
%%%   · 矢量矩阵用 \vect{} \mat{}，不要用 \mathbf 或 \boldsymbol
%%%   · 每个公式都给 \label，引用时用 \cref
%%% ---------------------------------------------------------------------------
\section{方法}
\label{sec:method}

\subsection{总体结构}
\label{sec:method:overview}

设骨干网络输出的层级特征为 $\setof{\featmap_l}_{l=2}^{5}$，其中 $\featmap_l \in
\R^{C \times H_l \times W_l}$，$H_l = H/2^{\,l}$。\ourmethod{} 在自顶向下
路径的每一次融合处替换原有的逐元素相加，代之以\cref{sec:method:gate} 定义的
门控算子。整体结构如\cref{fig:architecture} 所示。

\subsection{跨层注意力门控}
\label{sec:method:gate}

记上层特征经上采样对齐后为 $\tilde{\featmap}_{l+1}$。门控权重
$\mat{G}_l \in [0,1]^{H_l \times W_l}$ 由两层特征的通道统计量共同决定：

\begin{equation}
  \label{eq:gate}
  \mat{G}_l = \sigma\!\left(
    \mat{W}_2 \relu\!\left(
      \mat{W}_1 \left[ \operatorname{GAP}(\featmap_l) \,\Vert\,
                       \operatorname{GAP}(\tilde{\featmap}_{l+1}) \right]
    \right)
  \right),
\end{equation}

式中 $\sigma(\cdot)$ 为 Sigmoid 函数，$\operatorname{GAP}$ 为全局平均池化，
$\Vert$ 表示通道维拼接，$\mat{W}_1 \in \R^{C/r \times 2C}$ 与
$\mat{W}_2 \in \R^{C \times C/r}$ 为可学习参数，压缩比取 $r = 16$。

融合结果为

\begin{equation}
  \label{eq:fusion}
  \featmap_l' = \mat{G}_l \odot \featmap_l
              + (\mathbf{1} - \mat{G}_l) \odot \mathcal{D}(\tilde{\featmap}_{l+1}),
\end{equation}

其中 $\odot$ 为逐元素乘，$\mathcal{D}(\cdot)$ 为\cref{sec:method:align}
中的可变形对齐算子。由\cref{eq:fusion} 可见，当 $\mat{G}_l \to \mathbf{1}$ 时
融合退化为只保留浅层特征，当 $\mat{G}_l \to \mathbf{0}$ 时退化为标准
特征金字塔，因此本文方法是后者的严格推广。

\subsection{可变形对齐}
\label{sec:method:align}

上采样引入的偏移量 $\Delta \vect{p}$ 由 $3 \times 3$ 卷积从拼接特征回归：

\begin{equation}
  \label{eq:offset}
  \mathcal{D}(\tilde{\featmap})(\vect{p})
    = \sum_{k=1}^{9} w_k \cdot
      \tilde{\featmap}\!\left(\vect{p} + \vect{p}_k + \Delta \vect{p}_k\right),
\end{equation}

$\vect{p}_k$ 为标准卷积核的采样网格，$w_k$ 为核权重。

\subsection{理论分析}
\label{sec:method:theory}

下面说明门控融合不会放大特征的方差，这是训练稳定性的必要条件。

\begin{assumption}
  \label{assum:indep}
  $\featmap_l$ 与 $\mathcal{D}(\tilde{\featmap}_{l+1})$ 在各空间位置上
  互不相关，且方差分别为 $\sigma_1^2$ 与 $\sigma_2^2$。
\end{assumption}

\begin{theorem}[方差不增]
  \label{thm:variance}
  在\cref{assum:indep} 下，对任意 $\mat{G}_l \in [0,1]^{H_l \times W_l}$，
  由\cref{eq:fusion} 给出的融合特征满足
  \begin{equation}
    \label{eq:var-bound}
    \Var\!\left[\featmap_l'\right] \le \max\setof{\sigma_1^2, \sigma_2^2}.
  \end{equation}
\end{theorem}

\begin{proof}
  记 $g \in [0,1]$ 为 $\mat{G}_l$ 在某一位置的取值。由独立性，
  \begin{equation}
    \label{eq:var-expand}
    \Var\!\left[\featmap_l'\right]
      = g^2 \sigma_1^2 + (1-g)^2 \sigma_2^2 .
  \end{equation}
  右端是关于 $g$ 的二次函数，在 $[0,1]$ 上的最大值只可能在端点取得，
  分别为 $\sigma_1^2$ 与 $\sigma_2^2$，即得\cref{eq:var-bound}。
\end{proof}

\begin{remark}
  \label{rmk:relax}
  \cref{assum:indep} 的独立性在实践中并不严格成立，但实验中未观察到
  训练发散，说明该界在相关性较弱时仍具指导意义。
\end{remark}

\subsection{训练流程}
\label{sec:method:train}

完整的训练流程见\cref{alg:training}。

\begin{algorithm}[htbp]
  \caption{\ourmethod{} 的训练流程}
  \label{alg:training}
  \begin{algorithmic}[1]
    \Require 训练集 $\mathcal{S}$，骨干网络参数 $\params\subtext{bb}$，
             学习率 $\eta$，总轮数 $T$
    \Ensure  收敛后的模型参数 $\params^{*}$
    \State 用 ImageNet 预训练权重初始化 $\params\subtext{bb}$，其余参数用
           Kaiming 初始化
    \For{$t \gets 1$ \textbf{to} $T$}
      \ForAll{小批量 $\mathcal{B} \subset \mathcal{S}$}
        \State 前向计算层级特征 $\setof{\featmap_l}$
        \For{$l \gets 4$ \textbf{down to} $2$}
          \State 按\cref{eq:gate} 计算门控权重 $\mat{G}_l$
          \State 按\cref{eq:fusion} 得到融合特征 $\featmap_l'$
        \EndFor
        \State 计算总损失 $\loss = \loss\subtext{cls} + \lambda \loss\subtext{reg}$
        \State $\params \gets \params - \eta \nabla_{\params} \loss$
      \EndFor
      \State 按余弦策略更新 $\eta$
    \EndFor
    \State \Return $\params^{*} \gets \argmin_{\params} \loss$
  \end{algorithmic}
\end{algorithm}

% !TeX root = ../main.tex
%%% ---------------------------------------------------------------------------
%%% 实验
%%%
%%% 这一节演示图、子图、三线表、表注的规范写法：
%%%   · 表题在表**上方**，图题在图**下方**（跨规范一致要求）
%%%   · 表格只用 \toprule \midrule \bottomrule，不画竖线
%%%   · 数字列用 siunitx 的 S 列，按小数点对齐
%%%   · 坐标轴与表头标注写「量 / 单位」，如 t / s（GB/T 3101 推荐斜杠形式）
%%%   · 图必须是矢量格式（PDF/EPS），位图不低于 600 dpi
%%%
%%% 图片放在 figures/ 下，\includegraphics 时不用写目录名和扩展名。
%%% 下面的 \placeholder 是占位框，换成 \includegraphics 即可。
%%% ---------------------------------------------------------------------------
\section{实验}
\label{sec:exp}

\subsection{数据集与实现细节}
\label{sec:exp:setup}

实验在 \dataset{} 上进行，该数据集包含 \num{11268} 幅遥感图像、\num{1793658}
个标注实例，覆盖 18 个类别。按官方划分取 \qty{50}{\percent} 为训练集、
\qty{17}{\percent} 为验证集、\qty{33}{\percent} 为测试集。

原始大幅面图像切分为 $1024 \times 1024$ 的子图，重叠 \qty{200}{\pixel}。
骨干网络为 ResNet-50，优化器为 SGD，动量 \num{0.9}，权重衰减
\num{1e-4}，初始学习率 \num{0.01}，批大小 16，共训练 \qty{36}{\epoch}。
所有实验在 4 张 NVIDIA A100 上完成。

\subsection{总体结构}
\label{sec:exp:arch}

\ourmethod{} 的整体结构如\cref{fig:architecture} 所示。

\begin{figure}[htbp]
  \centering
  \placeholder{0.85\linewidth}{4.5cm}{figures/architecture.pdf}
  \caption{\ourmethod{} 的整体结构}
  \label{fig:architecture}
\end{figure}

\subsection{与现有方法的比较}
\label{sec:exp:comparison}

\cref{tab:main} 给出了与代表性方法的比较。所提方法在总体 mAP 与小目标
子集上均取得最佳结果，且推理速度未出现明显下降。

\begin{table}[htbp]
  \centering
  \caption{在 \dataset{} 测试集上与现有方法的比较}
  \label{tab:main}
  \begin{threeparttable}
    \begin{tabular}{
      l
      S[table-format=2.1]
      S[table-format=2.1]
      S[table-format=2.1]
      S[table-format=2.0]
      S[table-format=3.1]
    }
      \toprule
      {方法} & {mAP / \%} & {AP\textsubscript{S} / \%} &
      {AP\textsubscript{L} / \%} & {速度 / FPS} & {参数量 / M} \\
      \midrule
      Faster R-CNN\tnote{a} & 68.2 & 41.3 & 79.6 & 24 & 41.4 \\
      RetinaNet             & 69.7 & 43.0 & 80.1 & 28 & 37.9 \\
      FPN + ATSS            & 74.7 & 48.5 & 84.2 & 33 & 38.5 \\
      \midrule
      \textbf{\ourmethod{}（本文）} & \bfseries 78.4 & \bfseries 54.7 &
      \bfseries 85.9 & 31 & 42.1 \\
      \bottomrule
    \end{tabular}
    \begin{tablenotes}[flushleft]
      \footnotesize
      \item[a] 所有对比方法均使用 ResNet-50 骨干网络，在相同的数据划分与
        训练配置下复现。
      \item AP\textsubscript{S} 与 AP\textsubscript{L} 分别为面积小于
        $32^2$ 与大于 $96^2$ 像素的目标子集上的平均精度。
    \end{tablenotes}
  \end{threeparttable}
\end{table}

\subsection{消融实验}
\label{sec:exp:ablation}

为分离两个模块的贡献，在验证集上做了消融实验，结果见\cref{tab:ablation}。
可以看出，注意力门控贡献了 \qty{2.3}{\percent} 的 mAP 提升，可变形对齐
在此基础上再贡献 \qty{1.4}{\percent}，两者叠加没有出现相互抵消。

\begin{table}[htbp]
  \centering
  \caption{各模块的消融实验结果}
  \label{tab:ablation}
  \begin{tabular}{cc S[table-format=2.1] S[table-format=2.1]}
    \toprule
    {注意力门控} & {可变形对齐} & {mAP / \%} & {AP\textsubscript{S} / \%} \\
    \midrule
             &          & 74.7 & 48.5 \\
    $\checkmark$ &        & 77.0 & 52.6 \\
             & $\checkmark$ & 76.1 & 50.8 \\
    $\checkmark$ & $\checkmark$ & \bfseries 78.4 & \bfseries 54.7 \\
    \bottomrule
  \end{tabular}
\end{table}

\subsection{定性结果}
\label{sec:exp:qualitative}

\cref{fig:qualitative} 给出了两组典型场景的检测结果。在密集停放的车辆
场景（\cref{fig:qual-a}）中，基线方法出现了明显的漏检，而所提方法保持了
较完整的召回。

\begin{figure}[htbp]
  \centering
  \begin{subfigure}[b]{0.47\linewidth}
    \centering
    \placeholder{\linewidth}{3.6cm}{figures/qual-vehicle.pdf}
    \caption{密集车辆场景}
    \label{fig:qual-a}
  \end{subfigure}
  \hfill
  \begin{subfigure}[b]{0.47\linewidth}
    \centering
    \placeholder{\linewidth}{3.6cm}{figures/qual-ship.pdf}
    \caption{近岸船舶场景}
    \label{fig:qual-b}
  \end{subfigure}
  \caption{典型场景的检测结果对比}
  \label{fig:qualitative}
\end{figure}

% !TeX root = ../main.tex
%%% ---------------------------------------------------------------------------
%%% 结论
%%%
%%% 结论不是摘要的复述。GB/T 7713.2 要求结论「明确、精练、完整、准确」，
%%% 应当包含：
%%%   1. 本研究得到的结果说明了什么（要有具体数据支撑）
%%%   2. 与已有结论的异同
%%%   3. 本研究的局限
%%%   4. 尚待解决的问题与后续方向
%%% 不要在结论里引入新的实验数据或新的论点。
%%% ---------------------------------------------------------------------------
\section{结论}
\label{sec:conclusion}

针对遥感图像中小目标特征在深层网络中衰减的问题，本文提出了跨层注意力
门控与可变形对齐相结合的特征融合方法。理论分析表明，门控融合在特征
互不相关的假设下不会放大方差（\cref{thm:variance}），因而不损害训练
稳定性；\dataset{} 上的实验表明，所提方法将 mAP 从 \qty{74.7}{\percent}
提升至 \qty{78.4}{\percent}，小目标子集的平均精度提升
\qty{6.2}{\percent}，推理速度由 \qty{33}{\fps} 降至 \qty{31}{\fps}，
代价可以接受。

本研究存在两点局限。其一，\cref{assum:indep} 的独立性假设在相邻层级
特征高度相关时并不成立，此时\cref{eq:var-bound} 的界会变松；其二，
实验仅在光学遥感图像上验证，对 SAR 等其他成像模态的适用性尚未考察。

后续工作将从两方面展开：一是把门控权重的生成从通道统计量扩展到显式的
尺度先验，二是在多模态遥感数据上验证方法的迁移能力。


%%% ---- 结尾部分 ----
%%% 注意顺序：致谢在参考文献**之前**，附录在参考文献**之后**（GB/T 7713）
% !TeX root = ../main.tex
%%% ---------------------------------------------------------------------------
%%% 致谢
%%%
%%% 位置：参考文献**之前**（GB/T 7713）。
%%% 内容惯例：技术协助、数据/设备提供方、有益讨论者；
%%%           基金资助已在首页脚注列出，此处不必重复。
%%%
%%% 类选项 anonymous 开启时，整个 acknowledgements 环境会被自动丢弃，
%%% 正文一个字都不用改。注意 \end{acknowledgements} 必须**单独成行**，
%%% 这是 comment 宏包的限制。
%%% ---------------------------------------------------------------------------
\begin{acknowledgements}
感谢中国科学技术大学超算中心提供的计算资源支持；感谢匿名评审人提出的
建设性意见，使本文的理论分析部分得以完善。
\end{acknowledgements}


%%% lang=en 时把标题改成 {References}
\printbibliography[title = {参考文献}]

\appendix
% !TeX root = ../main.tex
%%% ---------------------------------------------------------------------------
%%% 附录
%%%
%%% 位置：参考文献**之后**（GB/T 7713）。main.tex 里的 \appendix 已经把
%%% 编号切成「附录 A」，图表公式自动变成 图 A1 / 表 A1 / 式 (A1)。
%%%
%%% 放什么：过长的推导、完整的参数表、补充数据——凡是放正文会打断阅读、
%%% 但审稿人可能需要核对的内容。
%%% ---------------------------------------------------------------------------

\section{方差上界的紧性讨论}
\label{app:tightness}

\cref{thm:variance} 给出的界在 $g \in \setof{0, 1}$ 时取到等号。对
$g \in (0,1)$，由\cref{eq:var-expand} 可以进一步给出

\begin{equation}
  \label{eq:tight}
  \Var\!\left[\featmap_l'\right]
    \ge \frac{\sigma_1^2 \sigma_2^2}{\sigma_1^2 + \sigma_2^2},
\end{equation}

等号在 $g^{*} = \sigma_2^2 / (\sigma_1^2 + \sigma_2^2)$ 处取得。这说明
门控不仅不放大方差，在中间取值处实际上起到了方差收缩的作用。

\section{完整超参数配置}
\label{app:hyperparams}

\begin{table}[htbp]
  \centering
  \caption{完整的训练超参数}
  \label{tab:hyperparams}
  \begin{tabular}{l l}
    \toprule
    {超参数} & {取值} \\
    \midrule
    优化器            & SGD（动量 \num{0.9}） \\
    初始学习率        & \num{0.01} \\
    学习率策略        & 余弦退火，预热 \num{500} 步 \\
    权重衰减          & \num{1e-4} \\
    批大小            & 16（每卡 4） \\
    训练轮数          & \qty{36}{\epoch} \\
    输入尺寸          & $1024 \times 1024$ \\
    数据增强          & 随机水平/垂直翻转、随机旋转 $90^{\circ}$ \\
    门控压缩比 $r$    & 16 \\
    损失权重 $\lambda$ & \num{1.0} \\
    \bottomrule
  \end{tabular}
\end{table}


\end{document}
